% File: abstract_kz.tex
% ReconX diploma --- Kazakh annotation (Аңдатпа)
% Note: pdfLaTeX renders Kazakh letters (Ә Ғ Қ Ң Ө Ұ Ү Һ І) via T2A fontenc
% inside a Russian babel environment; hyphenation falls back to Russian rules.
\selectlanguage{russian}
\chapter*{Аңдатпа}
\addcontentsline{toc}{chapter}{Аңдатпа}
\markboth{\foreignlanguage{russian}{АҢДАТПА}}{\foreignlanguage{russian}{АҢДАТПА}}
\begin{SingleSpace}
Берілген диплом жобасында веб-қосымшаларды сыртқы енуге сынау үшін \textbf{ReconX} фреймворкінің жобалануы, іске асырылуы және эксперименттік бағалануы ұсынылған. Ұсынылған фреймворк <<қара жәшік>> әдіснамасына сүйенеді және бастапқы коды мен есептік деректері тестілеушіге қолжетімсіз қашықтағы нысанға қатысты пассивті барлауды, белсенді сипаттарды анықтауды, конфигурацияны талдауды және осалдықтарды кең спектрде тексеруді біріктіреді.

Фреймворк \textbf{Python~3} тілінде іске асырылған және 15 орындау сатысына топтастырылған 46 мамандандырылған модульден тұратын конвейер ретінде ұйымдастырылған. Барлау модульдері қосалқы домендер, DNS-жазбалары, порттар және технологиялар туралы ақпарат жинайды; конфигурация модульдері \textbf{TLS}, қауіпсіздік тақырыптарын, \textbf{CORS}, аутентификациялық cookie-файлдар мен \textbf{JWT}-токендерді тексереді; белсенді модульдер \textbf{OWASP Top~10}-ның ең өзекті санаттарын тексереді, оның ішінде шағылысқан XSS, SQL-инъекциялар, IDOR, Host тақырыбының инъекциясы, кэшті улау, прототипті ластау, ашық қайта бағыттаулар, SSRF/SSTI, XXE, HTTP request smuggling және жарыс жағдайлары. Орталық \textit{табыстар тізілімі} әрбір нәтижені \textbf{CWE}, \textbf{CVSS}, OWASP және MITRE ATT\&CK метадеректері бар бірыңғай схемаға келтіреді; корреляттор мен активтер тәуекелі модулі ең осал активтерді басымдылыққа қояды. Жергілікті \textbf{үлкен тілдік модель} (Ollama / DeepSeek-R1) сканерлеу аяқталғаннан кейін түсіндірмелі талдамалық есепті қалыптастырады.

Эксперименттік бағалау екі толықтырушы жағдайда жүргізілді: әдейі осал жасалған \textit{OWASP Juice Shop} нысанына қарсы бақыланатын зертханалық іске қосу және өзгертетін сұрауларға жол бермейтін \texttt{bug\_bounty} режимінде орындалған \texttt{miras.app} академиялық басқару жүйесіне санкцияланған шынайы зерттеу. Фреймворк нақты нысандағы толық конвейерді шамамен 1\,мин~43\,с-та аяқтап, OWASP Top~10 санаттары бойынша 721 ерекше табыс берді, табылған 49 активті төрт тәуекел деңгейіне жіктеді (оның ішінде 5 \textit{сыни} деңгейдегі актив) және корреляттор арқылы екі синтетикалық крос-модульдік табыс шығарды. Анықталған осалдықтар екі сынақ нысанының күтілетін санаттарына сәйкес келді, жіберілген 4\,227 HTTP-сұраудың тек 1-уі ғана конфигурацияланған шектен тыс болды (және фреймворкпен жіберілмес бұрын бұғатталды), ал ЖИ арқылы құрылған атқарушылық қысқаша есеп ең жоғары тәуекелі бар активті есептің басына дұрыс шығарды.

Нәтижелер ұсынылған әдістің сыртқы веб-қосымшаларды енуге сынауды толық қолмен орындалатын тәсілмен салыстырғанда айтарлықтай қайта жаңартылатын, түсіндірмелі және уақыт жағынан тиімдірек ететінін растайды, бұл ретте операторға айқын профильдер жүйесі арқылы интрузивтік әрекеттерді бақылау мүмкіндігі сақталады (\textit{safe}, \textit{bug\_bounty}, \textit{deep}, \textit{intrusive}).

\par\noindent{\small\textbf{Кілт сөздер:} сыртқы енуге сыну; веб-қосымшалардың қауіпсіздігі; OWASP Top~10; барлау; осалдықтарды сканерлеу; табыстарды корреляциялау; ЖИ көмегімен қауіпсіздікті талдау; ReconX.}
\end{SingleSpace}
\selectlanguage{english}
\clearpage
